Las siguientes subsecciones documentan el comportamiento empírico de cada algoritmo
sobre un conjunto de cuatro archivos PDF de horarios generados por el portal de
Autoservicios BUAP. Los archivos procesados fueron:

\begin{itemize}
    \item \texttt{OmarHorarioChido.pdf} — 521{,}508 bytes
    \item \texttt{PauHorario.pdf} — 446{,}908 bytes
    \item \texttt{HorarioAlfredo.pdf} — 548{,}868 bytes
    \item \texttt{HorarioAri.pdf} — 453{,}448 bytes
\end{itemize}

Todos los archivos corresponden a horarios de una sola página, lo que produce
burst times pequeños y homogéneos —condición que influye directamente en el
comportamiento de Round Robin, como se analiza más adelante.

% ─────────────────────────────────────────────────────────────────
\subsubsection{FCFS (First Come, First Served)}

En FCFS los archivos se despachan en el orden exacto en que el sistema operativo
los entrega a través de \texttt{DirectoryStream}, sin ningún reordenamiento previo.
El proceso que llega primero ocupa la CPU hasta terminar; solo entonces el siguiente
puede iniciar.

\begin{table}[H]
\centering
\caption{Métricas de planificación — FCFS}
\label{tab:fcfs-metricas}
\begin{tabular}{|l|r|r|r|r|r|}
\hline
\textbf{Archivo} & \textbf{Llegada} & \textbf{Inicio} & \textbf{Fin} & \textbf{Espera} & \textbf{Turnaround} \\
\hline
OmarHorarioChido.pdf &   0 ms &   0 ms & 243 ms &   0 ms & 243 ms \\
PauHorario.pdf       &   0 ms & 244 ms & 270 ms & 244 ms & 270 ms \\
HorarioAlfredo.pdf   &   0 ms & 270 ms & 324 ms & 270 ms & 324 ms \\
HorarioAri.pdf       &   0 ms & 324 ms & 409 ms & 324 ms & 409 ms \\
\hline
\multicolumn{4}{|r|}{\textbf{Promedio}} & \textbf{209 ms} & \textbf{311 ms} \\
\hline
\end{tabular}
\end{table}

El primer archivo en llegar, \texttt{OmarHorarioChido.pdf}, no espera en absoluto
(espera = 0\,ms) y ocupa la CPU durante 243\,ms. Los tres archivos restantes
acumulan la espera del anterior, de modo que \texttt{HorarioAri.pdf} debe aguardar
324\,ms antes de ser atendido. Esta acumulación es la principal desventaja teórica
de FCFS: si el primer proceso en llegar es el más costoso, penaliza a todos los demás.
Sin embargo, en este caso los burst times son cortos y homogéneos, por lo que el
turnaround promedio de 311\,ms resulta el más bajo de los tres algoritmos evaluados.

% ─────────────────────────────────────────────────────────────────
\subsubsection{SJF (Shortest Job First)}

SJF reordena la cola antes de ejecutar, priorizando los archivos de menor tamaño
en bytes. El tamaño del archivo se utiliza como estimador del burst time: un PDF
más pequeño contiene menos contenido y se procesa más rápido.

El orden de llegada original era: \texttt{OmarHorarioChido.pdf},
\texttt{PauHorario.pdf}, \texttt{HorarioAlfredo.pdf}, \texttt{HorarioAri.pdf}.
Tras inspeccionar los tamaños, el scheduler reordenó la cola de la siguiente manera:

\begin{enumerate}
    \item \texttt{PauHorario.pdf} — 446{,}908 bytes
    \item \texttt{HorarioAri.pdf} — 453{,}448 bytes
    \item \texttt{OmarHorarioChido.pdf} — 521{,}508 bytes
    \item \texttt{HorarioAlfredo.pdf} — 548{,}868 bytes
\end{enumerate}

\begin{table}[H]
\centering
\caption{Métricas de planificación — SJF}
\label{tab:sjf-metricas}
\begin{tabular}{|l|r|r|r|r|r|}
\hline
\textbf{Archivo} & \textbf{Llegada} & \textbf{Inicio} & \textbf{Fin} & \textbf{Espera} & \textbf{Turnaround} \\
\hline
PauHorario.pdf       &   0 ms &   0 ms & 173 ms &   0 ms & 173 ms \\
HorarioAri.pdf       &   0 ms & 173 ms & 227 ms & 173 ms & 227 ms \\
OmarHorarioChido.pdf &   0 ms & 227 ms & 432 ms & 227 ms & 432 ms \\
HorarioAlfredo.pdf   &   0 ms & 433 ms & 853 ms & 433 ms & 853 ms \\
\hline
\multicolumn{4}{|r|}{\textbf{Promedio}} & \textbf{208 ms} & \textbf{421 ms} \\
\hline
\end{tabular}
\end{table}

La espera promedio de SJF (208\,ms) es marginalmente menor que la de FCFS (209\,ms),
lo que confirma el principio teórico: atender primero los trabajos más cortos reduce
la espera acumulada de los procesos que llegan después. Sin embargo, el turnaround
promedio (421\,ms) es considerablemente mayor que el de FCFS, porque
\texttt{HorarioAlfredo.pdf} —el más pesado— terminó en 853\,ms, elevando el promedio.
Esto ilustra que SJF favorece los procesos cortos a costa de los más largos,
que en casos extremos pueden sufrir inanición.

% ─────────────────────────────────────────────────────────────────
\subsubsection{RR (Round Robin, quantum = 150 ms)}

Round Robin asigna a cada proceso un quantum de tiempo de 150\,ms. Si un archivo
no termina dentro de ese límite, el hilo es interrumpido, el archivo se reencola
al final de la cola circular y se retoma en la siguiente vuelta.

Durante la primera vuelta, \texttt{OmarHorarioChido.pdf} y \texttt{PauHorario.pdf}
excedieron el quantum y fueron reencolados. \texttt{HorarioAlfredo.pdf} y
\texttt{HorarioAri.pdf} terminaron dentro del límite. En la segunda vuelta,
ambos archivos pendientes completaron su procesamiento.

\begin{table}[H]
\centering
\caption{Métricas de planificación — Round Robin (quantum = 150 ms)}
\label{tab:rr-metricas}
\begin{tabular}{|l|r|r|r|r|r|}
\hline
\textbf{Archivo} & \textbf{Llegada} & \textbf{Inicio} & \textbf{Fin} & \textbf{Espera} & \textbf{Turnaround} \\
\hline
OmarHorarioChido.pdf &   0 ms &   9 ms & 464 ms &   9 ms & 464 ms \\
PauHorario.pdf       &   0 ms & 161 ms & 477 ms & 161 ms & 477 ms \\
HorarioAlfredo.pdf   &   0 ms & 311 ms & 419 ms & 311 ms & 419 ms \\
HorarioAri.pdf       &   0 ms & 419 ms & 443 ms & 419 ms & 443 ms \\
\hline
\multicolumn{4}{|r|}{\textbf{Promedio}} & \textbf{225 ms} & \textbf{450 ms} \\
\hline
\end{tabular}
\end{table}

Round Robin produjo la espera promedio más alta (225\,ms) y el turnaround promedio
más alto (450\,ms) de los tres algoritmos. Esto se explica por la naturaleza de
los archivos: los PDFs de horarios BUAP son pequeños y homogéneos, con burst times
que rondan los 100--300\,ms. En este escenario, el overhead de interrumpir y
reencolar procesos supera el beneficio de compartir la CPU, y RR converge al
comportamiento de FCFS con penalización adicional por el cambio de contexto simulado.

% ─────────────────────────────────────────────────────────────────
\subsubsection*{Comparativa general}

\begin{table}[H]
\centering
\caption{Resumen comparativo de los tres algoritmos}
\label{tab:comparativa-algoritmos}
\begin{tabular}{|l|r|r|}
\hline
\textbf{Algoritmo} & \textbf{Espera promedio} & \textbf{Turnaround promedio} \\
\hline
FCFS                        & 209 ms & 311 ms \\
SJF                         & 208 ms & 421 ms \\
RR (quantum = 150 ms)       & 225 ms & 450 ms \\
\hline
\end{tabular}
\end{table}

Para cargas de trabajo homogéneas y de burst time corto —como los horarios de una
página generados por Autoservicios BUAP— FCFS resulta el algoritmo más eficiente
en términos de turnaround. SJF ofrece la menor espera promedio teórica pero penaliza
a los procesos más pesados. Round Robin, diseñado para cargas heterogéneas y sistemas
interactivos, introduce overhead innecesario cuando los procesos son uniformes,
resultando en el peor desempeño para este caso de uso particular.
